\chapter{Stacks Project dependency table}
\label{app:stacks}

\newcommand{\stacktag}[1]{%
  \href{https://stacks.math.columbia.edu/tag/#1}{Tag~#1}}

This appendix records the algebraic-geometry results from the Stacks Project
that are used in the book. Its purpose is not to replace the Stacks Project,
but to prevent a stable tag from being cited as though it said more than it
does. For each tag we give three pieces of information:
\begin{enumerate}
\item the exact statement imported into the book;
\item the hypotheses that must be checked at the point of use;
\item the first application in the logical order of the manuscript.
\end{enumerate}
The online statements were checked on 14 August 2026, with Tag 0892 added
and rechecked on 25 August 2026. Tags are stable
identifiers, whereas section and lemma numbers can change as the Stacks
Project is reorganized. All citations below refer to
\cite{stacks-project}.

\begin{remark}[Direct and inherited dependencies]
\label{rem:stacks-direct-inherited}
A tag is a \emph{direct dependency} if a proof in this book invokes its
conclusion. It is an \emph{inherited dependency} if we import a theorem
whose proof in the Stacks Project invokes that tag and the intermediate step
is relevant to our hypothesis audit. We record both kinds, but never count
an inherited dependency as an additional theorem proved here.
\end{remark}

\begin{remark}[The scheme boundary]
\label{rem:stacks-scheme-boundary}
Every result in this appendix is stated for schemes or rings. When a chapter
concerns Berkovich analytifications, the tag supplies only the algebraic
input. Topological claims about analytifications---compactness,
Hausdorffness, openness, or extension of seminorms---are proved separately
in the relevant chapter. In particular, the algebraic valuative criterion
for properness must not be used as an analytic valuative criterion for
compact complex spaces.
\end{remark}

\section{Valuations, centers, and flatness}
\label{sec:stacks-valuations}

\subsection*{\stacktag{01KZ}: separated morphisms have valuative uniqueness}

\index{valuative criterion!separatedness}
The statement is:
\begin{quote}
If \(f:X\to S\) is a separated morphism of schemes, then \(f\) satisfies the
uniqueness part of the valuative criterion.
\end{quote}
Thus two morphisms \(\Spec R\to X\), where \(R\) is a valuation ring, that
agree on \(\Spec\operatorname{Frac}(R)\) and have the same composite to
\(S\), are equal.
There is no finite-type or noetherian hypothesis in this direction.

The first application is \Cref{prop:centers-proper-separated}.
Take \(S=\Spec k\), \(R=R_v\), and use the generic-point map
\(\Spec F\to X\). Separatedness gives uniqueness of the center. The proof
in \cref{chap:valuations} also gives the equalizer argument directly; the tag records the
general theorem containing that argument.

\subsection*{\stacktag{0BX5}: the valuative criterion for properness}

\index{valuative criterion!properness}
Let \(S\) be a scheme and let \(f:X\to Y\) be a morphism of schemes over
\(S\). Under the hypotheses that \(f\) is of finite type and
quasi-separated, the following are equivalent:
\begin{enumerate}
\item \(f\) is proper;
\item \(f\) satisfies the valuative criterion;
\item for every valuation ring \(R\) with fraction field \(K\), every solid
diagram
\[
\begin{tikzcd}
\Spec K \arrow[r] \arrow[d] & X \arrow[d,"f"]\\
\Spec R \arrow[r] \arrow[ru,dashed] & Y
\end{tikzcd}
\]
has a unique dotted arrow.
\end{enumerate}
The finite-type and quasi-separated assumptions belong to the equivalence
and may not be dropped when using the converse direction.

For \Cref{prop:centers-proper-separated}, only
\[
                         \text{proper}\Longrightarrow
                         \text{existence and uniqueness}
\]
is used. A projective variety is proper, finite type, and separated over
the ground field, so the hypotheses are automatic there.

\subsection*{\stacktag{01KV}: cancellation for separatedness}

\index{separated morphism!cancellation}
For composable morphisms of schemes
\[
                         X\xrightarrow{f}Y\xrightarrow{g}Z,
\]
if \(g\circ f\) is separated, then \(f\) is separated. The same assertion
holds with “quasi-separated” in place of “separated.” This tag is
\emph{not} the valuative uniqueness theorem; that role is played by
\stacktag{01KZ}.

Tag 01KV is an inherited dependency of the Noetherian Stein factorization
in \stacktag{03H0}: it is used there to see that the connected-fibre factor
is separated, hence proper after the remaining finite-type and universal
closedness conditions have been established. The distinction explains the
corrected citation in \cref{chap:valuations}.

\subsection*{\stacktag{0539}: torsion-free modules over a valuation ring}

\index{valuation ring!torsion-free module}
If \(A\) is a valuation ring and \(M\) is an \(A\)-module, then
\[
                         M\text{ is flat over }A
                         \quad\Longleftrightarrow\quad
                         M\text{ is torsion-free}.
\]
No finite-generation assumption is present.

The first application is \Cref{lem:torsion-free-flat-valuation}.
For the schematic closure of a generic-fibre subscheme, one works affinely
and obtains an injection
\[
              A/I\lhook\joinrel\longrightarrow
              (A\otimes_{K^\circ}K)/I_K.
\]
The target is a vector space over \(K\), hence torsion-free over
\(K^\circ\). Therefore the quotient \(A/I\) is torsion-free and Tag 0539
makes the closure flat. Notice that the argument needs the displayed
injection; flatness does not follow merely from the word “closure.”

\section{Blowups, flattening, and normalization}
\label{sec:stacks-flattening}

\subsection*{\stacktag{081T}: domination by an admissible blowup}

\index{blowup!admissible}
Let \(\varphi:X\to S\) be proper, with \(S\) quasi-compact and
quasi-separated. Let \(U\subset S\) be a quasi-compact open over which
\(\varphi\) is an isomorphism. Then there is a \(U\)-admissible blowup
\[
                         S'\longrightarrow S
\]
and a factorization \(S'\to X\to S\).

Here \(U\)-admissible means that the center is disjoint from \(U\).
The conclusion says that a blowup of the target dominates the given proper
modification; it does not flatten a sheaf. The first explicit pointer is
the discussion following \Cref{thm:flattening-models}. In the
envelope argument it is used only when a proper modification is already
known to be an isomorphism over the chosen dense open.

\subsection*{\stacktag{0815}: flattening by blowup}

\index{flattening by blowup}
Let \(S\) be quasi-compact and quasi-separated, let \(X\) be a
quasi-compact \(S\)-scheme locally of finite presentation, and let
\(\mathcal F\) be a finite-type quasi-coherent \(\mathcal O_X\)-module.
Let \(U\subset S\) be quasi-compact. If \(\mathcal F_U\) is finitely
presented and flat over \(U\), then a \(U\)-admissible blowup
\(S'\to S\) makes the strict transform \(\mathcal F'\) finitely presented
and flat over \(S'\).

The first application is \Cref{thm:flattening-models}. There the
base is noetherian, the morphism is of finite presentation, and
\(\mathcal F\) is coherent. These stronger assumptions imply the
quasi-compactness and finite-presentation conditions in the tag. The
second major application is \Cref{lem:bj-finite-flatification}, with
\[
                         \mathcal F=h_*\mathcal O_T
\]
and \(U\) its dense flat locus. The algebra structure survives passage to
the strict transform because the torsion subsheaf being removed is a
two-sided ideal.

\begin{remark}[Two different blowup theorems]
\label{rem:stacks-two-blowup-theorems}
Tags 081T and 0815 are often cited together, but their outputs are different:
\[
\begin{array}{c@{\qquad}c}
\stacktag{081T}&\text{a blowup dominating a proper modification},\\
\stacktag{0815}&\text{a blowup flattening a specified module}.
\end{array}
\]
A proof using the second output must name the module and its strict
transform. A proof using the first must name the proper modification and
the open set over which it is an isomorphism.
\end{remark}

\subsection*{\stacktag{02K9}: finite locally free morphisms}

\index{finite locally free morphism}
This is a section tag. The particular result used here is Lemma 29.49.2:
\[
\begin{split}
f\text{ finite locally free}
&\Longleftrightarrow
f\text{ finite, flat, and locally of finite presentation}.
\end{split}
\]
If the base is locally noetherian, “finite and flat” already implies
finite locally free. The section also records stability under composition
and base change and the open-and-closed decomposition by rank.

In \Cref{lem:bj-finite-flatification}, the strict transform
\(\mathcal A_0\) is finite, flat, and finitely presented. Tag 02K9
therefore turns
\[
                  \Spec_{Z_0}(\mathcal A_0)\longrightarrow Z_0
\]
into a finite locally free morphism of constant generic rank
\(d=[K(T):K(X)]\). Faithful flatness additionally uses that the rank is
positive on the integral connected base.

\subsection*{\stacktag{01OF}: the blowup section}

\index{blowup!integrality}
This is the Stacks Project section on blowups. The precise item used in
finite flatification is Lemma 31.33.9:
\begin{quote}
If \(X\) is integral and \(\mathcal I\subset\mathcal O_X\) is a nonzero
quasi-coherent ideal, then the blowup of \(X\) in \(\mathcal I\) is
integral.
\end{quote}
The section also identifies the blowup with the relative Proj of the Rees
algebra, proves that it is an isomorphism away from the center, and proves
projectivity when the ideal is of finite type.

The first application is again \Cref{lem:bj-finite-flatification}.
The admissible center misses the dense open \(U\), and the blowup is
birational over the integral target. Equivalently, its Rees algebra embeds
in \(K(X)[t]\), which is the elementary reason integrality is preserved.
The tag is not a statement that an arbitrary blowup of a reducible scheme
is integral.

\subsection*{\stacktag{0892}: ample twists}

Let $f:X\to Y$ be a morphism of schemes, let $\cM$ be an invertible
$\cO_Y$-module, and let $\cL$ be an invertible $\cO_X$-module.  The part
used in this book says that if $\cL$ is $f$-ample and $\cM$ is ample, then
\[
                     \cL\otimes f^*\cM^{\otimes a}
\]
is ample for all sufficiently large integers $a$.  The hypotheses imply the
quasi-compactness needed for one exponent to work globally.

The first application is
\Cref{thm:p01-real-pl-positive-current-candidate}.  Take $f=\mathrm{id}_V$
on the projective residual model, $\cM=A$, and
$\cL=M_i$ or $M_i^{-1}$.  Every invertible sheaf is relatively ample for
the identity morphism, so the tag gives that
$A^{\otimes N}\otimes M_i^{\pm1}$ is ample for all large $N$.  Finiteness
of the family permits one common $N$.

\subsection*{\stacktag{035S}: finite normalization over a Nagata scheme}

\index{normalization!finiteness}
If \(X\) is a Nagata scheme, its normalization
\[
                         \nu:X^\nu\longrightarrow X
\]
is finite. The Nagata hypothesis is essential to the cited statement.
Every finite-type scheme over a field is excellent and hence Nagata, so the
tag applies to all projective varieties in \cref{part:bj}.

The first application in finite descent is the normalization of the
finite-type integral blowup \(Z_0\) in
\Cref{lem:bj-finite-flatification}. It is also used in
\Cref{cor:bj-normalization-source-ep} and
\Cref{lem:bj-unibranch-normalization-universal}. Finiteness then
implies projectivity over a projective target, but projectivity is not part
of Tag 035S itself.

\section{Finite group actions and determinant norms}
\label{sec:stacks-finite}

\subsection*{\stacktag{0BRI}: transitivity on primes}

\index{Galois action!primes}
Let a finite group \(G\) act on a ring \(R\). Any two primes of \(R\)
lying over the same prime of the invariant ring \(R^G\) are conjugate under
\(G\).

The first application is the separable part of
\Cref{lem:bj-finite-normal-open}. For the Galois normalization
\(S\to X\), this tag shows that the support primes of analytic lifts of a
fixed point of \(X^{\An}\) form one \(G\)-orbit. It says nothing yet about
the absolute values on the residue fields; that is why Tag 0BRJ and a
separate valuation-extension argument are also needed.

\subsection*{\stacktag{0BRJ}: decomposition groups and residue fields}

\index{decomposition group}
In the setting of Tag 0BRI, fix a prime \(\mathfrak q\subset R\) above
\(\mathfrak p\subset R^G\), and let
\[
                     D=\{\sigma\in G\mid\sigma(\mathfrak q)=\mathfrak q\}.
\]
Then \(\kappa(\mathfrak q)/\kappa(\mathfrak p)\) is algebraic and normal,
and the natural map
\[
       D\longrightarrow
       \Aut\bigl(\kappa(\mathfrak q)/\kappa(\mathfrak p)\bigr)
\]
is surjective.

In \Cref{lem:bj-finite-normal-open}, this gives the group action on
the normal residue extension. Transitivity on extensions of an absolute
value is then a theorem about valuation rings in a finite normal field
extension, not an extra conclusion of Tag 0BRJ. A purely inseparable part
has a unique extension and must be treated separately.

\subsection*{\stacktag{0BD2}: determinant norms of functions}

\index{norm!finite locally free}
If \(\pi:X\to Y\) is finite locally free of degree \(d\geq1\), multiplication
by \(f\in\pi_*\mathcal O_X\) is an endomorphism of a rank-\(d\) vector
bundle. Its determinant defines a canonical norm
\[
                         \operatorname{Norm}_\pi(f)=\det(m_f)
\]
of degree \(d\), and this construction commutes with arbitrary base change.

The first application is the universal norm polynomial in
\Cref{chap:bj-finite-descent}. After adjoining variables \(T_j\),
the determinant of multiplication by \(\sum_jT_js_j\) produces the
coefficients used in the weighted fibre average. Tag 0BD2 directly covers
functions after trivializing a line bundle. Multiplicativity for line
bundles and the global determinant norm line
\[
       \det(q_*M)\otimes\det(q_*\mathcal O_W)^{-1}
\]
are taken from Ferrand's norm functor \cite{FerrandNorm}; they are not
attributed to Tag 0BD2 alone.

\section{Stein factorization and relative normalization}
\label{sec:stacks-stein}

\subsection*{\stacktag{03H0}: Noetherian Stein factorization}

\index{Stein factorization}
Let \(S\) be locally noetherian and let \(f:X\to S\) be proper. There is a
factorization
\[
                         X\xrightarrow{f'}S'\xrightarrow{\pi}S
\]
such that:
\begin{enumerate}
\item \(f'\) is proper with geometrically connected fibres;
\item \(\pi\) is finite;
\item \(f'_*\mathcal O_X=\mathcal O_{S'}\);
\item \(S'=\Spec_S(f_*\mathcal O_X)\);
\item \(S'\) is the normalization of \(S\) in \(X\).
\end{enumerate}
The phrase “normalization of \(S\) in \(X\)” is relative; it is not in
general the normalization of \(S\) in its own total ring of fractions.

The first application is the factorization
\[
                         Y\xrightarrow{g}T\xrightarrow{h}X
\]
of a smooth projective alteration in
\Cref{chap:bj-alterations}. Projective varieties over a field are
noetherian and the alteration is proper, so the hypotheses hold. Tag 03H0
gives connected fibres and finiteness; birationality of \(g\) is deduced
separately from generic finiteness and the identity \(K(T)=K(Y)\).

\subsection*{\stacktag{035L}: normality of relative normalization}

\index{relative normalization}
Let \(f:Y\to X\) be quasi-compact and quasi-separated, and let \(X'\) be
the normalization of \(X\) in \(Y\). Assume:
\begin{enumerate}
\item \(Y\) is normal;
\item every quasi-compact open of \(Y\) has finitely many irreducible
components.
\end{enumerate}
Then \(X'\) is a disjoint union of integral normal schemes, and
\(Y\to X'\) is dominant and induces a bijection on irreducible components.

In \Cref{chap:bj-alterations}, the source \(Y\) is a smooth integral
projective variety. It is noetherian, so its quasi-compact opens have
finitely many irreducible components. The Stein middle term \(T\) is the
relative normalization by Tag 03H0. Tag 035L therefore makes \(T\) normal;
because \(Y\) has one irreducible component, it also makes \(T\) integral.
Finiteness of \(T\to X\) comes from Tag 03H0, not from Tag 035L.

\section{Universal homeomorphisms and subintegral extensions}
\label{sec:stacks-universal-homeomorphisms}

\subsection*{\stacktag{01S4}: radicial equals universally injective}

\index{radicial morphism}
For a morphism \(f:X\to S\), the following conditions are equivalent:
\begin{enumerate}
\item maps \(\Spec K\to X\) over \(S\) are unique for every field \(K\);
\item \(f\) is universally injective;
\item \(f\) is radicial;
\item the diagonal \(X\to X\times_SX\) is surjective.
\end{enumerate}
Radicial means injective on points with purely inseparable residue-field
extensions.

The first application is the definition-and-criterion discussion in
\Cref{chap:bj-universal-homeomorphisms}. Later, in
\Cref{lem:bj-unibranch-normalization-universal}, bijectivity of the
normalization together with purely inseparable residue extensions gives
radiciality. Surjectivity is checked separately.

\subsection*{\stacktag{04DF}: criterion for a universal homeomorphism}

\index{universal homeomorphism}
A morphism \(f:X\to Y\) is a universal homeomorphism if and only if it is
\[
                 \text{integral},\qquad
                 \text{universally injective},\qquad
                 \text{surjective}.
\]
The forward direction includes integrality, not merely universal
bijectivity. The reverse direction uses that integral morphisms are
universally closed.

In \Cref{lem:bj-unibranch-normalization-universal}, the normalization
is finite, hence integral; Tag 01S4 gives universal injectivity; and
normalization is surjective. Tag 04DF then yields a universal
homeomorphism. The induced homeomorphism on analytifications requires the
additional uniqueness of absolute-value extensions through purely
inseparable residue fields.

\subsection*{\stacktag{0CNF}: a characteristic-\(p\) affine criterion}

\index{universal homeomorphism!characteristic p criterion}
Let \(p\) be prime and let \(A\to B\) induce an isomorphism
\(A[1/p]\to B[1/p]\). The map
\(\Spec B\to\Spec A\) is a universal homeomorphism if and only if:
\begin{enumerate}
\item the kernel is locally nilpotent;
\item for every \(b\in B\), there is a \(p\)-power \(q\) such that both
\(qb\) and \(b^q\) lie in the image of \(A\to B\).
\end{enumerate}
When \(A\subset B\) are algebras of characteristic \(p\), the first term
\(qb\) is zero and the useful conclusion is \(b^q\in A\).

The first application is \Cref{lem:bj-uniform-frobenius}. Since a
finite algebra has finitely many generators on each member of a finite
affine cover, finitely many elementwise \(p\)-powers can be replaced by one
common power. The Frobenius identity then handles every polynomial in the
generators. The tag itself gives only an exponent depending on \(b\); the
uniform exponent is a separate finite-generation argument.

\subsection*{\stacktag{0CND}: square--cube towers}

\index{subintegral extension!square--cube criterion}
For an injective ring extension \(A\subset B\), the following are equivalent:
\begin{enumerate}
\item \(\Spec B\to\Spec A\) is a universal homeomorphism inducing
isomorphisms on all residue fields;
\item every finite subset of \(B\) is contained in an intermediate algebra
\[
                         A[b_1,\ldots,b_n]\subset B
\]
such that
\[
             b_i^2,b_i^3\in A[b_1,\ldots,b_{i-1}]
             \qquad(1\leq i\leq n).
\]
\end{enumerate}
This is the constructive form of subintegrality used in characteristic
zero.

The first application is \Cref{chap:bj-subintegral-normalization}.
It reduces finite pieces of a subintegral normalization to a succession of
elementary square--cube extensions. The later additive-power argument uses
the stronger Roberts--Singh theorem; Tag 0CND supplies the structural
subintegrality and its finite-stage presentation.

\subsection*{\stacktag{0EUH}: base change of square--cube extensions}

\index{subintegral extension!base change}
Let \(\varphi:R\to S\) have locally nilpotent kernel and suppose \(S\) is
generated over \(R\) by elements \(x\) satisfying
\[
                         x^2,x^3\in\varphi(R).
\]
Then \(\Spec S\to\Spec R\) is a homeomorphism and all residue-field maps are
isomorphisms. After any ring map \(R\to R'\), the base-changed map
\[
                         R'\longrightarrow R'\otimes_RS
\]
again has locally nilpotent kernel and is generated by elements satisfying
the same square--cube condition.

Together with Tag 0CND, this proves that polynomial base change preserves
the subintegral extensions used in
\Cref{lem:bj-local-additive-sandwich}. In particular,
\[
                         A\subset B
             \quad\Longrightarrow\quad
                         A[\mathbf u,\mathbf z]
                         \subset B[\mathbf u,\mathbf z]
\]
remains subintegral. This step is needed before applying the
Roberts--Singh additive-power theorem to a generic polynomial. Tag 0EUH
does not supply that additive-power theorem.

\section{Dependency chains used in the envelope theorem}
\label{sec:stacks-dependency-chains}

The long proof in \cref{part:bj} uses the preceding results in four short chains.
Writing them explicitly makes the hypotheses visible.

\subsection*{Finite normal descent}

For a finite dominant map \(T\to X\) between normal projective varieties,
the algebraic chain is
\[
\begin{aligned}
\text{Galois normalization}
 &\xrightarrow[\text{Nagata}]{\stacktag{035S}}
 \text{finite cover},\\
\text{group action}
 &\xrightarrow{\stacktag{0BRI},\,\stacktag{0BRJ}}
 \text{prime and residue control},\\
h_*\mathcal O_T
 &\xrightarrow{\stacktag{0815}}
 \text{flat transform}
 \xrightarrow{\stacktag{02K9}}
 \text{finite locally free},\\
\text{finite locally free}
 &\xrightarrow{\stacktag{0BD2}}
 \text{coefficient norms}.
\end{aligned}
\]
The first line enters analytic openness; the last two enter the
multiplicity-weighted fibre average. Neither chain alone proves descent.

\subsection*{Alterations}

For a smooth alteration \(Y\to X\) with \(X\) normal,
\[
  Y\to X
  \xrightarrow{\stacktag{03H0}}
  \bigl(Y\to T\to X\bigr)
  \xrightarrow{\stacktag{035L}}
  T\text{ normal and integral}.
\]
The first arrow in the factorization is then shown birational by a
function-field calculation, while the second is finite by Tag 03H0.
Birational reflection and finite normal descent, proved in separate
chapters, may therefore be applied in that order.

\subsection*{Positive-characteristic normalization}

For the normalization of a unibranch variety,
\[
\begin{aligned}
\text{unibranch normalization}
 &\xrightarrow{\stacktag{01S4},\,\stacktag{04DF}}
 \text{universal homeomorphism},\\
\text{affine characteristic-\(p\) piece}
 &\xrightarrow{\stacktag{0CNF}}
 \text{elementwise powers},\\
\text{finite generators}
 &\Longrightarrow
 \text{one exponent on the cover}.
\end{aligned}
\]
The last implication is an argument in \cref{chap:bj-universal-homeomorphisms},
not part of the tag.

\subsection*{Characteristic-zero normalization}

For a finite subintegral extension of \(\mathbf Q\)-algebras,
\[
\begin{aligned}
\text{subintegral}
 &\xrightarrow{\stacktag{0CND}}
 \text{finite square--cube towers},\\
\text{square--cube tower}
 &\xrightarrow{\stacktag{0EUH}}
 \text{base-change stability}.
\end{aligned}
\]
Roberts--Singh then supplies the additive-power formula used in the ideal
sandwich. The separation of these steps is important: the Stacks Project
tags provide the algebraic form of subintegrality, while the uniform
asymptotic identity is a different theorem.

\section{Compact audit table}
\label{sec:stacks-compact-table}

For quick reference, the following ledger gives the first place at which
each direct dependency is used. “Inherited” means
\Cref{rem:stacks-direct-inherited}.
\[
\begin{array}{lll}
\toprule
\text{Tag}&\text{imported content}&\text{first application}\\
\midrule
01KZ&\text{separated implies valuative uniqueness}&\text{\cref{chap:valuations}}\\
0BX5&\text{valuative criterion for properness}&\text{\cref{chap:valuations}}\\
01KV&\text{separatedness cancellation}&\text{inherited through 03H0}\\
0539&\text{torsion-free equals flat over a valuation ring}&\text{\cref{chap:models}}\\
081T&\text{domination by an admissible blowup}&\text{\cref{chap:models}}\\
0815&\text{flattening by admissible blowup}&\text{\cref{chap:models}}\\
035S&\text{finite normalization of a Nagata scheme}&\text{\cref{chap:bj-finite-descent}}\\
0BRI&\text{transitivity on primes under a finite group}&\text{\cref{chap:bj-finite-descent}}\\
0BRJ&\text{decomposition group on residue fields}&\text{\cref{chap:bj-finite-descent}}\\
02K9&\text{finite flat versus finite locally free}&\text{\cref{chap:bj-finite-descent}}\\
01OF&\text{integrality and projectivity of blowups}&\text{\cref{chap:bj-finite-descent}}\\
0892&\text{an ample twist dominates a fixed invertible sheaf}&\text{\cref{chap:delta-forms}}\\
0BD2&\text{base-change-compatible determinant norm}&\text{\cref{chap:bj-finite-descent}}\\
03H0&\text{Noetherian Stein factorization}&\text{\cref{chap:bj-alterations}}\\
035L&\text{normality of relative normalization}&\text{\cref{chap:bj-alterations}}\\
01S4&\text{radicial equals universally injective}&\text{\cref{chap:bj-universal-homeomorphisms}}\\
04DF&\text{criterion for universal homeomorphism}&\text{\cref{chap:bj-universal-homeomorphisms}}\\
0CNF&\text{characteristic-\(p\) affine criterion}&\text{\cref{chap:bj-universal-homeomorphisms}}\\
0CND&\text{square--cube tower criterion}&\text{\cref{chap:bj-subintegral-normalization}}\\
0EUH&\text{base change of square--cube extensions}&\text{\cref{chap:bj-subintegral-normalization}}\\
\bottomrule
\end{array}
\]

\begin{remark}[Scope of the ledger]
\label{rem:stacks-ledger-scope}
The table lists the stable tags on which proofs in the present first draft
materially depend. Routine background---for example, that finite morphisms
are affine, that projective morphisms are proper, or that finite-type
schemes over a field are noetherian---is not assigned a separate row.
Whenever such a routine fact becomes the delicate step of a later proof, it
should be promoted to the ledger with its precise tag.
\end{remark}

\begin{remark}[Source note]
\label{rem:stacks-source-note}
The tag selection for the envelope chapters follows the algebraic-tools
section of the supplied manuscript
\emph{The envelope property for unibranch projective varieties}, with two
clarifications made explicit here. First, Tag 01KZ, rather than Tag 01KV,
is the direct separatedness input for uniqueness of centers. Second,
Tags 081T and 0815 have distinct outputs and are never treated as a single
flattening theorem.
\end{remark}
